\chapter{Introduction}
\label{ch:introduction}

\section{Background to the Research}

The Standard Model has been tested to extraordinary precision, yet it is
not considered a complete theory because it says nothing about gravity at the quantum level,
it has no dark matter candidate despite dark matter making up roughly a
quarter of the universe's energy budget (Planck Collaboration, 2020), and
it cannot explain why the universe is full of matter and almost no
antimatter. These gaps motivate the search for physics beyond the Standard
Model (BSM), one well-motivated form of which is a class of light, weakly
coupled bosonic fields, such as axions, axion-like particles, and dark
photons, that mediate feeble, long-range \emph{fifth forces}.

This thesis is concerned with a specific subset of fifth forces: those that
couple to fermion spin rather than to charge or mass alone. Spin has no
classical analogue, which makes \emph{exotic spin-dependent interactions}
an unusually clean channel for probing symmetry violation. Two languages
are used to describe such interactions. The Standard Model Extension (SME)
of Colladay and Kostelecký (1997, 1998) is a relativistic effective field
theory that systematically parametrises Lorentz and CPT violation through
fixed background coefficients such as $b_\mu$, $H_{\mu\nu}$, and
$d_{\mu\nu}$. The Dobrescu--Mocioiu (DM) classification (Dobrescu and
Mocioiu, 2006) is a complete, non-relativistic basis of sixteen potentials,
$V_1$--$V_{16}$, in which nearly all laboratory searches for exotic forces
report their bounds. Connecting the two requires the Foldy--Wouthuysen (FW)
transformation (Foldy and Wouthuysen, 1950), the standard technique for
extracting the non-relativistic limit of a relativistic Dirac Hamiltonian.

Separately, antimatter physics has become one of the sharpest available
tools for testing fundamental symmetries. Experiments at CERN, such as BASE
and ALPHA, compare matter and antimatter magnetic moments and spectral
lines at a precision that directly tests charge-parity-time (CPT)
invariance (Smorra \textit{et al.}, 2017; Ahmadi \textit{et al.}, 2017).
Combining these two threads, exotic spin-dependent interaction searches and
antimatter precision physics, opens a regime in which a CPT-violating
interaction would appear as a measurable asymmetry between the coupling
constants measured for matter and for antimatter.

\section{Research Problem and Hypotheses}

Despite over a hundred experiments compiled through 2024 (Cong \textit{et
al.}, 2025), no unified framework systematically translates between the
SME and DM formalisms, and no framework exists for comparing matter-sector
and antimatter-sector constraints on exotic spin-dependent interactions
within a single, quantitative measure of CPT consistency. The SME describes
Lorentz- and CPT-violating physics through operator coefficients that are
not directly expressed as measurable potentials; the DM basis is
experimentally convenient but does not make its relativistic origin
explicit. Prior attempts at a mapping between the two are partial,
restricted to subsets of coefficients or potentials, and no compiled
constraint database exists in a form that allows a direct, standardised
comparison of matter- and antimatter-sector bounds across particle species,
interaction types, and experimental platforms.

This research problem yields two working hypotheses examined in this
thesis. First, that a complete Foldy--Wouthuysen reduction of the four SME
coefficients that generate spin-dependent terms, the CPT-odd $b_\mu$ and
$g_{\lambda\mu\nu}$ and the CPT-even $H_{\mu\nu}$ and $d_{\mu\nu}$, can be
matched, coefficient by coefficient, onto the DM potential basis, and that
the resulting potentials will change sign between matter and antimatter for
the CPT-odd coefficients but not for the CPT-even ones. Second, that a
dimensionless asymmetry parameter $A_\alpha$, computed from currently
published experimental bounds, can be used to test this prediction
quantitatively once matter- and antimatter-sector datasets are compiled
and standardised.

The first hypothesis is confirmed for all four coefficients, with one
clarification about what is being compared. The sign flip for the CPT-odd
$b_\mu$ and $g_{\lambda\mu\nu}$, and its absence for the CPT-even
$H_{\mu\nu}$ and $d_{\mu\nu}$, hold between CPT-conjugate states, in which
the antiparticle's spin is reversed. At the same physical spin the pattern
is reversed, as Kostelecký and Lane (1999b) state
(Section~\ref{sec:results-fw-bmu}). The second
hypothesis is limited by the nature of the data. Published bounds are
one-sided upper limits rather than signed measurements, so $A_\alpha$ alone
cannot always separate a genuine CPT-violating signal from an ordinary gap
in experimental sensitivity (Chapters~3 and~4). This thesis treats that
limitation as a central finding rather than an incidental caveat.

\section{Justification for the Research}

A unified SME--DM mapping gives experimentalists a direct way to translate
a measured or constrained potential into a bound on a relativistic,
symmetry-classified coefficient, and vice versa, rather than relying on
partial, coefficient-specific translations scattered across the
literature. A standardised, cross-platform constraint database removes the
unit and convention inconsistencies that currently make it difficult to
compare bounds from nitrogen-vacancy-centre magnetometry, torsion balances,
atomic co-magnetometers, positronium spectroscopy, and antimatter
spectroscopy on equal footing. Together, these tools enable the first
systematic matter--antimatter comparison of spin-dependent forces and
provide a reusable, extensible pipeline (SPINDEP) rather than a one-off
calculation, one that automatically incorporates future antimatter-sector
measurements as they are published. The results are relevant to ongoing
experimental programmes at CERN (BASE, ALPHA), to precision quantum sensing
laboratories, and to the SME data tables maintained by Kostelecký and
Russell (2011), which this work is intended to complement.

\section{Aims and Objectives}

The primary objective of this research is to develop a unified theoretical
and computational framework that connects Standard Model Extension
coefficients to Dobrescu--Mocioiu spin-dependent potentials and enables
systematic comparison of matter- and antimatter-sector constraints to test
CPT symmetry. The specific objectives are:

\begin{enumerate}
  \item to derive explicit mappings between the SME fermion-sector
    coefficients $b_\mu$, $H_{\mu\nu}$, $d_{\mu\nu}$, and $g_{\lambda\mu\nu}$
    and the Dobrescu--Mocioiu coupling structures $g_s$, $g_p$, $g_V$, and $g_A$
    using the Foldy--Wouthuysen transformation;
  \item to construct a unified, standardised database of experimental
    constraints from nitrogen-vacancy-centre magnetometry, torsion
    balances, atomic co-magnetometers, positronium spectroscopy, and
    antimatter experiments;
  \item to define and compute the CPT asymmetry parameter $A_\alpha$ across
    every matter--antimatter pair the compiled database can support;
  \item to perform a global coverage analysis identifying unconstrained or
    weakly constrained regions of the parameter space, particularly in the
    antimatter sector; and
  \item to propose concrete experimental and data-compilation strategies
    for closing the identified gaps.
\end{enumerate}

\section{Outline of the Thesis}

\begin{itemize}
  \item Chapter~1 introduces the motivation, research questions, and
    objectives of the thesis.
  \item Chapter~2 reviews the theoretical background, the experimental
    literature, and the specific inadequacies in existing knowledge that
    this thesis addresses.
  \item Chapter~3 describes the materials and methods: the
    Foldy--Wouthuysen procedure used to derive the SME--DM mapping, and
    the database compilation and statistical methodology used to compute
    the CPT asymmetry parameter.
  \item Chapter~4 presents the results, the derived coefficient-to-potential
    mapping, the compiled 283-dataset constraint database, the
    matter--antimatter asymmetry analysis, and the coverage gap analysis,
    together with their discussion.
  \item Chapter~5 presents the conclusions of the thesis and outlines its
    principal limitations and directions for future work.
\end{itemize}
